\exercisesfor{5}

The conventions of § 5 remain valid unless otherwise mentioned.

\begin{enumerate}
\def\labelenumi{\arabic{enumi}.}
\item
  Let g be the 2-dimensional non-commutative solvable Lie algebra. Show that the Killing form of g is non-zero, that every invariant bilinear form on g is degenerate and that every derivation of g is inner.
\item
  \begin{enumerate}
  \def\labelenumii{(\alph{enumii})}
  \tightlist
  \item
    Show that in the 3-dimensional solvable Lie algebra over \(\mathbf{R}\) defined by the multiplication table \([x, y] = z\) , \([x, z] = -y\) , \([y, z] = 0\) , there exists no decreasing sequence of ideals of dimensions 3, 2, 1, 0.
  \end{enumerate}
\end{enumerate}

\begin{enumerate}
\def\labelenumi{(\alph{enumi})}
\setcounter{enumi}{1}
\tightlist
\item
  Show that in the non-commutative 2-dimensional solvable Lie algebra g there exists a sequence of ideals of dimensions 2, 1, 0 but that g is not nil-potent.
\end{enumerate}

\begin{enumerate}
\def\labelenumi{\arabic{enumi}.}
\setcounter{enumi}{2}
\item
  Let \(\mathfrak{g}\) be a solvable Lie algebra such that the conditions \(x\in \mathfrak{g},y\in \mathfrak{g}\) , \([[x,y],y] = 0\) imply \([x,y] = 0\) . Show that \(\mathfrak{g}\) is commutative. (Let \(k\) be the largest integer such that \(\mathcal{D}^{k - 1}\mathfrak{g}\neq \{0\}\) , \(\mathcal{D}^k\mathfrak{g} = \{0\}\) . Assuming \(k\geqslant 2\) , show first that \([\mathcal{D}^{k - 2}\mathfrak{g},\mathcal{D}^{k - 1}\mathfrak{g}] = \{0\}\) and then that \([\mathcal{D}^{k - 2}\mathfrak{g},\mathcal{D}^{k - 2}\mathfrak{g}] = \{0\}\) , whence a contradiction.)
\item
  Show that the centre of \(\mathfrak{st}(n,\mathbf{K})\) is zero and that that of \(\mathfrak{n}(n,\mathbf{K})\) is of dimension 1.
\item
  Let \(\mathfrak{g}\) be a Lie algebra and \((\mathcal{D}^0\mathfrak{g},\mathcal{D}^1\mathfrak{g},\ldots ,\mathcal{D}^n\mathfrak{g})\) the sequence of derived algebras of \(\mathfrak{g}\) ( \(n\geqslant 0,\mathcal{D}^{n - 1}\mathfrak{g}\neq \mathcal{D}^n\mathfrak{g}\) ). Then \(\dim \mathcal{D}^i\mathfrak{g} / \mathcal{D}^{i + 1}\mathfrak{g}\geqslant 2^{i - 1} + 1\) for \(1\leqslant i\leqslant n - 2\) . (Taking quotients by \(\mathcal{D}^n\mathfrak{g}\) , reduce it to the case where \(\mathfrak{g}\) is solvable. Then use the fact that \(\mathcal{D}_{\mathfrak{g}}\) is nilpotent and Exercise 7 (c) of § 4.)
\item
  \begin{enumerate}
  \def\labelenumii{(\alph{enumii})}
  \tightlist
  \item
    Verify that the following multiplication table defines a 5-dimensional solvable Lie algebra g:
  \end{enumerate}
\end{enumerate}

\[
\begin{array}{r} \left[ x _ {1}, x _ {2} \right] = x _ {5}, \qquad \left[ x _ {1}, x _ {3} \right] = x _ {3}, \qquad \left[ x _ {2}, x _ {4} \right] = x _ {4} \\ \left[ x _ {1}, x _ {4} \right] = \left[ x _ {2}, x _ {3} \right] = \left[ x _ {3}, x _ {4} \right] = \left[ x _ {5}, \mathfrak {g} \right] = 0. \end{array}
\]

\begin{enumerate}
\def\labelenumi{(\alph{enumi})}
\setcounter{enumi}{1}
\item
  Show that the orthogonal of \(\mathfrak{g}\) with respect to the Killing form is \(\mathcal{D}\mathfrak{g} = \mathrm{K}x_3 + \mathrm{K}x_4 + \mathrm{K}x_5\) . Deduce that \(\mathcal{D}\mathfrak{g}\) is the largest nilpotent ideal of \(\mathfrak{g}\) .
\item
  Show that there exists no supplementary subalgebra of \(\mathcal{D}\mathfrak{g}\) in \(\mathfrak{g}\) . Deduce that \(\mathfrak{g}\) is not the semi-direct product of a commutative algebra and a nilpotent ideal. (Show that this nilpotent ideal would necessarily be \(\mathcal{D}\mathfrak{g}\) .)
\end{enumerate}

\begin{enumerate}
\def\labelenumi{\arabic{enumi}.}
\setcounter{enumi}{6}
\item
  Let \(\mathfrak{g}\) be the 3-dimensional solvable Lie algebra with basis \((x,y,z)\) such that \([x,y] = z\) , \([x,z] = y\) , \([y,z] = 0\) . Show that the linear mapping which maps \(x\) to \(-x, y\) to \(-z, z\) to \(y\) is an automorphism of \(\mathfrak{g}\) of order 4. Compare this result with Exercise 21 (c) of §4.
\item
  \begin{enumerate}
  \def\labelenumii{(\alph{enumii})}
  \tightlist
  \item
    Let \(\mathfrak{g}_0\) be a 3-dimensional solvable real Lie algebra such that \(\mathcal{D}\mathfrak{g}_0\) is commutative and 2-dimensional. Let \(\mathfrak{g}\) be the algebra derived from \(\mathfrak{g}_0\) by extending the base field from \(\mathbf{R}\) to \(\mathbf{C}\) . For \(x \in \mathfrak{g}\) , let \(u_x\) be the restriction of \(\mathrm{ad}_{\mathfrak{g}}x\) to \(\mathcal{D}\mathfrak{g}\) . Show that the eigenvalues of \(u_x\) either have the same absolute value or are linearly dependent over \(\mathbf{R}\) . (We have \(x = \lambda y + z\) with \(z \in \mathcal{D}\mathfrak{g}\) , \(y \in \mathfrak{g}_0\) , \(\lambda \in \mathbf{C}\) and hence \(u_x = \lambda u_y\) . Now \(u_y\) is the \(\mathbf{C}\) -linear extension to \(\mathcal{D}\mathfrak{g}\) of an \(\mathbf{R}\) -linear endomorphism of \(\mathcal{D}\mathfrak{g}_{0}\) .)
  \end{enumerate}
\end{enumerate}

\begin{enumerate}
\def\labelenumi{(\alph{enumi})}
\setcounter{enumi}{1}
\item
  Show that there exists a 3-dimensional solvable complex Lie algebra g, with \(\mathcal{D}_{\mathfrak{g}}\) commutative and 2-dimensional, and an element \(x\) of g such that the restriction of \(\mathrm{ad}_{\mathfrak{g}}x\) to \(\mathcal{D}_{\mathfrak{g}}\) has eigenvalues which neither have the same absolute value nor are linearly dependent over R. (Construct g as a semi-direct product of a 1-dimensional algebra and a 2-dimensional commutative algebra.)
\item
  Show that the algebra constructed in (b) cannot be derived from a real Lie algebra by extending the base field from \(\mathbf{R}\) to \(\mathbf{C}\) .
\end{enumerate}

\begin{enumerate}
\def\labelenumi{\arabic{enumi}.}
\setcounter{enumi}{8}
\tightlist
\item
  Let \(\mathfrak{g}\) be a Lie algebra, \(\mathfrak{r}\) its radical, \(\mathfrak{n}\) its largest nilpotent ideal and D a derivation of \(\mathfrak{g}\) . Show that \(\mathcal{D}\mathfrak{g} \cap \mathfrak{r} \subset \mathfrak{n}\) . (Let \(\mathfrak{d}\) be the Lie algebra of derivations of \(\mathfrak{g}\) and \(\mathfrak{r}'\) its radical. If \(x \in \mathfrak{g}\) is such that \(\mathrm{D}x \in \mathfrak{r}\) , then
\end{enumerate}

\[
\operatorname{ad} (\mathrm{D} x) = [ \mathrm{D}, \operatorname{ad} x ] \in \mathscr {D} \mathfrak {d} \cap \mathfrak {r} ^ {\prime}
\]

(Corollary 2 to Proposition 5), hence \(\operatorname{ad}(\mathbf{D}x)\) is nilpotent (Theorem 1) and hence \(\mathbf{D}x \in \mathfrak{n}\) .)

\begin{enumerate}
\def\labelenumi{\arabic{enumi}.}
\setcounter{enumi}{9}
\item
  Let g be a Lie algebra, r its radical and a a subinvariant subalgebra of g. Show that the radical of a is \(a \cap r\) . (Apply Corollary 3 to Proposition 5 several times.)
\item
  Let \(\mathfrak{g}\) be a Lie algebra, \(\rho\) a finite-dimensional representation of \(\mathfrak{g}\) and \(\mathbb{E}\) the associative endomorphism algebra generated by 1 and \(\rho(\mathfrak{g})\) . Show that the largest nilpotency ideal \(n\) of \(\rho\) is equal to the set \(n'\) of \(x \in \mathfrak{g}\) such that \(\operatorname{Tr}(\rho(x)u) = 0\) for all \(u \in E\) . (To show that \(n' \subset n\) , show that \(n'\) is an ideal and that for all \(x \in n'\) the semi-simple component of \(\rho(x)\) is zero, by noting that \(\operatorname{Tr}((\rho(x))^n) = 0\) for every integer \(n > 0\) .)
\item
  Suppose that K is algebraically closed. Let g be a nilpotent Lie K-algebra. Let \(\rho\) be a finite-dimensional representation of g on a vector space V. For every linear form \(\lambda\) on g, let \(V^{\lambda}\) be the vector subspace of V consisting of the \(\xi \in V\) such that, for all \(x \in g\) , \((\rho(x) - \lambda(x)\mathrm{I})^{n}\xi = 0\) for sufficiently large n.
\end{enumerate}

\begin{enumerate}
\def\labelenumi{(\alph{enumi})}
\item
  The \(V^{\lambda}\) are stable under \(\rho(\mathfrak{g})\) and the sum of the \(V^{\lambda}\) is direct. (Use Exercise 22 of § 4.)
\item
  We have \(\mathrm{V} = \sum \mathrm{V}^{\lambda}\) . (If each \(\rho(x)\) has a single eigenvalue, \(\mathrm{V} = \mathrm{V}^{\lambda_0}\) by Corollary 2 to Theorem 1. If \(\rho(x_0)\) has at least two distinct eigenvalues, \(\mathrm{V}\) )
\end{enumerate}

is the direct sum of two non-trivial subspaces which are stable under \(\rho(\mathfrak{g})\) . Then argue by induction on the dimension of V.)

\begin{enumerate}
\def\labelenumi{\arabic{enumi}.}
\setcounter{enumi}{12}
\item
  Let \(\mathfrak{g}\) be a Lie Algebra and \(\mathfrak{D}\) the Lie algebra of derivations of \(\mathfrak{g}\) . For \(\mathfrak{g}\) to be characteristically nilpotent, it is necessary and sufficient that \(\mathfrak{D}\) be nilpotent and that \(\dim \mathfrak{g} > 1\) . (To see that the condition is sufficient, write \(\mathfrak{g}\) as a sum of subspaces \(\mathfrak{g}^{\lambda}\) , applying Exercise 12 to the identity mapping of \(\mathfrak{D}\) . Show that \([\mathfrak{g}^{\lambda}, \mathfrak{g}^{\mu}] \subset \mathfrak{g}^{\lambda + \mu}\) and that each \(\mathfrak{g}^{\lambda}\) is an ideal of \(\mathfrak{g}\) . Deduce that \(\mathfrak{g}^{\lambda}\) is commutative for \(\lambda \neq 0\) . Using again the fact that \(\mathfrak{D}\) is nilpotent, show that \(\mathfrak{g} = \mathfrak{g}^{0}\) if \(\dim \mathfrak{g} > 1\) . Otherwise \(\mathfrak{g} = \mathfrak{g}^{0} \times \mathfrak{h}\) , where \(\mathfrak{h}\) is commutative. Show first that \(\dim \mathfrak{h} \leqslant 1\) . If \(\dim \mathfrak{h} = 1\) , note that there would exist a derivation D of \(\mathfrak{g}\) such that D( \(\mathfrak{g}^{0}\) ) = \{0\} and D( \(\mathfrak{h}\) ) is contained in the centre of \(\mathfrak{g}^{0}\) ( \(\S 4\) , Exercise 8 ( \(a\) )).)
\item
  Let V be a finite-dimensional vector space over K and z an endomorphism of V. We adopt the notation \(z_{q}^{p}\) of Exercise 3 of §3. An endomorphism \(z'\) of V is called a replica of z if for all p and q every zero of \(z_{q}^{p}\) is a zero of \(z_{q}^{\prime p}\) . Show that if \(\operatorname{Tr}(zz') = 0\) for every replica \(z'\) of \(z\) , then z is nilpotent. (Use the proof of Lemma 3. In the notation of this proof, prove in particular that t is a replica of z.)
\item
  Let K be a field of characteristic 2. The identity representation of the nilpotent Lie algebra \(\mathfrak{sl}(2,\mathbf{K})\) on \(K^{2}\) defines a semi-direct product h of \(\mathfrak{sl}(2,\mathbf{K})\) by \(K^{2}\) . Show that h is solvable but that Dh is not nilpotent. Deduce that h admits no faithful linear representation by triangular matrices. Show also that the conclusions of Exercise 5 are false.
\item
  Let g be a Lie algebra over an arbitrary field K, A a commutative associative algebra over K and \(g' = g \otimes_{K} A\) , which can be considered as a Lie algebra over K.
\end{enumerate}

\begin{enumerate}
\def\labelenumi{(\alph{enumi})}
\item
  If D is a derivation of A, show that there exists one and only one derivation \(D'\) of \(g'\) such that \(\mathrm{D}'(x \otimes a) = x \otimes \mathrm{D}a\) for \(x \in g, a \in A\) .
\item
  Let \(p\) be a prime number, \(G\) a cyclic group of order \(p\) and \(s\) a generator of \(G\) . Suppose that \(K\) is of characteristic \(p\) and henceforth let \(A\) be the algebra of \(G\) over \(K\) . Show that there exists one and only one derivation \(D\) of \(A\) such that \(D(s^k) = ks^{k-1}\) for all \(k \in \mathbf{Z}\) . Show that the \(K\) -linear combinations of the \(x \otimes (s - 1)^k\) ( \(k = 1, 2, \ldots, p - 1, x \in \mathfrak{g}\) ) form a solvable ideal \(\tau\) of \(g'\) and that \(g'/\tau\) is isomorphic to \(\mathfrak{g}\) .
\item
  Take \(\mathfrak{g}\) to be simple (cf.~§6, no. 2, Definition 2). Then \(\mathfrak{r}\) is the radical of \(\mathfrak{g}'\) but is not a characteristic ideal. (Observe that \(\mathrm{D}(x\otimes (s - 1)) = 1\otimes x.\) )
\end{enumerate}

\begin{enumerate}
\def\labelenumi{\arabic{enumi}.}
\setcounter{enumi}{16}
\tightlist
\item
  Let g be a Lie algebra. Suppose that for every simple g-module M of finite dimension over K the \(x_{M}\) are permutable with one another. Show that g is solvable. (Observe that Dg is contained in the nilpotent radical and hence is solvable.)
\end{enumerate}

